\documentclass[10pt,a4paper,sans]{moderncv}
\moderncvstyle{classic}
\moderncvcolor{blue}
\usepackage[utf8]{inputenc}
\usepackage[scale=0.82]{geometry}

% moderncv compose \today en anglais par défaut ; format français sans charger babel.
\renewcommand{\today}{\number\day\space\ifcase\month\or janvier\or février\or mars\or avril\or
mai\or juin\or juillet\or août\or septembre\or octobre\or novembre\or décembre\fi\space\number\year}

\name{Alexandro}{Disla}
\title{Lettre de Motivation}
\address{Rue Saint-Louis Jeanty \#14, Route de Frère}{Port-au-Prince, Haïti}{}
\phone[mobile]{(+509) 4148-3700}
\email{alexandrodisla@hotmail.com}
\extrainfo{GitHub: \url{https://github.com/AD0791}}

\begin{document}

\recipient{Équipe de Recrutement --- Catholic Relief Services Haïti}{Port-au-Prince, Haïti\\Poste : Senior Project Officer MEAL --- Projet RAPID}
\date{\today}
\opening{Madame, Monsieur,}
\closing{Je vous prie d'agréer, Madame, Monsieur, l'expression de mes salutations distinguées,}
\enclosure[Ci-joint]{Curriculum Vitae détaillé, copies de diplômes}

\makelettertitle

Je vous adresse ma candidature au poste de Senior Project Officer MEAL du projet RAPID en Haïti. J'apporte cinq ans d'expérience terrain en suivi, évaluation, redevabilité et apprentissage sur des projets multi-sites haïtiens, répartis sur trois secteurs : la santé et le VIH à la Caris Foundation International, l'eau et l'assainissement chez HANWASH, l'éducation chez Anseye Pou Ayiti. Le français et le créole haïtien sont mes langues maternelles, je rédige mes rapports en anglais comme en français, et je suis disponible pour Port-au-Prince comme pour Les Cayes ou Fort-Liberté, avec les déplacements que le poste demande.

Ce que vous décrivez sous le terme d'ICT4D est le cœur de ma pratique. J'ai conçu des instruments de collecte sur ODK, KoboToolbox, CommCare et mWater, avec logiques de saut et contraintes de validation, afin que les erreurs soient arrêtées à la saisie plutôt que corrigées trois semaines plus tard ; j'ai administré les bases relationnelles qui les reçoivent et construit les tableaux de bord Power BI, Quarto et mWater qui les rendent lisibles. À la Caris Foundation, j'ai automatisé la production des rapports d'activité en Python et SQL, réduisant de 40 \% le temps de traitement manuel --- du temps rendu à l'analyse plutôt qu'à la mise en forme. Peu de profils MEAL administrent eux-mêmes le système qu'ils pilotent ; c'est mon autre métier.

Sur la qualité des données et la conformité bailleur, mon expérience porte précisément sur ce qui doit résister à une vérification. J'ai établi et supervisé les protocoles d'assurance qualité d'un projet de santé multi-sites, en remontant les chiffres rapportés jusqu'aux registres sources et en suivant les actions correctives jusqu'à leur clôture. Je produisais et soumettais les indicateurs MER destinés à l'USAID via DATIM : le rythme, les contrôles de cohérence et les échéances d'un bailleur du gouvernement américain me sont familiers, ce qui devrait servir un projet financé par le Département d'État. Sur la restitution, j'ai formé des équipes de terrain aux outils et à la qualité des données, et rédigé les guides techniques qui leur permettent de continuer sans moi --- la matière même des événements trimestriels de réflexion que vous décrivez, où l'analyse doit revenir aux équipes sous une forme qui change ce qu'elles font.

Un point que je préfère poser franchement. À la Caris Foundation, j'ai dirigé techniquement les agents de collecte de l'ensemble des sites pendant près de trois ans --- répartition du travail, contrôle qualité, formation aux procédures --- et servi de référent technique au personnel de suivi-évaluation. En revanche, les évaluations de performance formelles, les plans de développement individuel et le recrutement ne relevaient pas de mon mandat : la gestion hiérarchique au sens où CRS la décrit serait pour moi un premier exercice. J'ai la pratique de la conduite d'équipe sur le terrain, il me manque le cadre RH qui l'accompagne, et je préfère l'annoncer plutôt que de le laisser découvrir en entretien. De même, ma pratique du MEAL s'est faite en santé, en EAH et en éducation plutôt qu'en réponse d'urgence ; l'exigence de rigueur, elle, ne change pas de secteur.

Ce poste m'attire parce que l'approche convergente de RAPID --- les mêmes ménages suivis à travers plusieurs services essentiels --- pose exactement le problème de données que je sais résoudre : identifier un ménage de façon fiable d'un secteur à l'autre, sans le compter deux fois ni le perdre. Je serais heureux de mettre cela au service d'une équipe qui travaille auprès des populations les plus vulnérables du pays, et je reste à votre entière disposition pour un entretien.

\makeletterclosing

\end{document}
